\chapter{DSL generico per i vincoli}
\label{ch:dsl}
\label{ch:dsl_generico_per_i_vincoli}

Linguaggio compatto, dichiarativo, totalmente componibile per
esprimere vincoli HARD/SOFT/PREFERRED/ENFORCED su qualunque
combinazione logica di predicati sul modello dei dati.
Implementato in \texttt{webui/backend/utils/general\_dsl.py}
(parser ricorsivo discendente + AST + evaluator post-hoc, $\sim$250
LOC, no dipendenze esterne). Esposto via
\texttt{POST /api/constraints/general}, validato live da
\texttt{POST /api/constraints/general/validate}, accessibile dal
modal ``+ Nuovo vincolo DSL'' del tab Vincoli.

\section{Filosofia}

Prima del DSL generico il sistema aveva quattro sub-DSL
specializzati (query, logical, objective, introspettivo). Stessa
grammatica implementata 4 volte. Il DSL generico unifica tutto
sotto una sola grammatica con quattro ``sapori'' di compilazione
(LIST, LOGIC, OBJECTIVE, EVAL): \emph{un parser, molti
compilatori}. La sintassi di
\texttt{forall x in S where Filter: Body} e dei predicati atomici
(\texttt{==}, \texttt{!=}, \texttt{<}, \texttt{in}, \dots) e'
identica in tutti i contesti; cambia solo il backend di
traduzione.

Il DSL e' interpretato \emph{post-hoc}: parsea l'espressione e
la valuta DOPO che la pipeline (Phase A, Phase B, metaeuristiche)
ha prodotto una soluzione candidata. Le violazioni HARD sono
rilevate dal Feasibility Check; le SOFT alimentano lo score
globale.

\section{Grammatica BNF}

\begin{lstlisting}[basicstyle=\ttfamily\small]
expr           ::= iff_expr
iff_expr       ::= implies_expr ( "<=>" implies_expr )*
implies_expr   ::= or_expr     ( "=>"  or_expr )*
or_expr        ::= and_expr    ( ("OR"|"or") and_expr )*
and_expr       ::= not_expr    ( ("AND"|"and") not_expr )*
not_expr       ::= ("NOT"|"not") not_expr | atom

atom           ::= quant_expr | count_expr | comparison
                 | call | bool_lit | "(" expr ")"

quant_expr     ::= ("forall"|"exists") IDENT "in" source
                   ["where" or_expr] ":" expr

count_expr     ::= "count" IDENT "in" source ["where" or_expr]
                   ( ":" comparison | op value )

source         ::= IDENT                  -- literal name
                 | IDENT ("." IDENT)+     -- path (runtime)

comparison     ::= value ( op value
                         | "in" "[" value ("," value)* "]"
                         | "in" path
                         | "not_in" "[" value ("," value)* "]"
                         | "not_in" path )

op             ::= "==" | "!=" | "<" | "<=" | ">" | ">="
value          ::= IDENT ( "." IDENT )* | NUM | STRING
                 | BOOL | function_call
function_call  ::= IDENT "(" [arg_list] ")"
arg_list       ::= arg ("," arg)*
arg            ::= IDENT "==" value | value
\end{lstlisting}

Precedenze (dal pi\`u basso al pi\`u alto):
\texttt{<=>}, \texttt{=>}, \texttt{OR}, \texttt{AND}, \texttt{NOT},
confronti, \texttt{.} (member access). Le parentesi sovrascrivono.

\section{Sorgenti iterabili}

\begin{tabular}{ll}
\toprule
\textbf{Sorgente} & \textbf{Tipo elemento} \\
\midrule
\texttt{lessons}     & lezioni della soluzione attiva \\
\texttt{assignments} & cattedre (Assignment rows) \\
\texttt{teachers}    & docenti (\texttt{name, group, max\_hours, subject}) \\
\texttt{classes}     & classi (\texttt{name, year, curriculum, n\_students}) \\
\texttt{classrooms}  & aule (\texttt{name, kind, type, tags}) \\
\texttt{students}    & studenti (con \texttt{tags[]} e \texttt{groups[]}) \\
\texttt{subjects}    & materie \\
\texttt{curricula}   & indirizzi di studio \\
\texttt{groups}      & gruppi articolati \\
\texttt{slots}       & le 36 \texttt{(day, hour)} della settimana \\
\texttt{days}        & i 6 giorni (Lun..Sab; index 1..6) \\
\texttt{hours}       & le 6 ore (8..13) \\
\bottomrule
\end{tabular}

\textbf{Sorgente-path}: dalla v0.5 puoi scrivere
\texttt{exists g in s.groups: g.name == "X"} dove \texttt{s.groups}
e' un attributo che restituisce una lista; il parser accetta
qualunque path puntato dopo \texttt{in}.

\section{Galleria di esempi}

Ogni esempio sotto e' stato validato live contro il parser
corrente; copia-incolla nell'editor del modal ``+ Nuovo vincolo
DSL'' e funziona.

\subsection{Vincoli docente (10 esempi)}

\begin{enumerate}
\item Borghi non insegna mai in 1A alla 6a ora:
\begin{lstlisting}
forall l in lessons where l.teacher == Borghi
                       and l.class == 1A:
    l.hour != 13
\end{lstlisting}

\item Borghi: max 4 ore/giorno totali:
\begin{lstlisting}
forall d in days:
    count l in lessons where l.teacher == Borghi
                          and l.day == d.index: l <= 4
\end{lstlisting}

\item Lab fisica: ogni docente di Fisica esattamente 1 ora a
settimana in lab\_fisica:
\begin{lstlisting}
forall t in teachers where t.subject == "Fisica":
    count l in lessons where l.teacher == t
        and l.classroom.type == "lab_fisica": l == 1
\end{lstlisting}

\item Tetto cattedra Inglese: max 18h/settimana:
\begin{lstlisting}
forall t in teachers where t.subject == "Inglese":
    count l in lessons where l.teacher == t: l <= 18
\end{lstlisting}

\item Rossi non lavora il sabato (HARD):
\begin{lstlisting}
forall l in lessons where l.teacher == Rossi
                       and l.day == 6:
    false
\end{lstlisting}

\item Rossi: al pi\`u 1 sesta ora/settimana:
\begin{lstlisting}
count l in lessons where l.teacher == Rossi
                      and l.hour == 13: l <= 1
\end{lstlisting}

\item Tutti A050: max 5 ore/giorno:
\begin{lstlisting}
forall t in teachers where t.group == "A050":
    forall d in days:
        count l in lessons where l.teacher == t
            and l.day == d.index: l <= 5
\end{lstlisting}

\item Coupling Rossi 3A $\Rightarrow$ Bianchi 3B:
\begin{lstlisting}
(exists l in lessons:
    l.teacher == Rossi and l.class == 3A)
=> (exists l in lessons:
    l.teacher == Bianchi and l.class == 3B)
\end{lstlisting}

\item Rossi mai in palestra:
\begin{lstlisting}
forall l in lessons where l.teacher == Rossi:
    l.classroom.type != "palestra"
\end{lstlisting}

\item Tutti i prof di mate: almeno una lezione di mattina:
\begin{lstlisting}
forall t in teachers where t.subject == "Matematica":
    exists l in lessons where l.teacher == t:
        l.hour <= 9
\end{lstlisting}
\end{enumerate}

\subsection{Vincoli classe (5 esempi)}

\begin{enumerate}
\item La 1A non lavora il sabato:
\begin{lstlisting}
forall l in lessons where l.class == 1A: l.day != 6
\end{lstlisting}

\item Le prime: max 5 ore/giorno:
\begin{lstlisting}
forall c in classes where c.year == 1:
    forall d in days:
        count l in lessons where l.class == c.name
            and l.day == d.index: l <= 5
\end{lstlisting}

\item Quarte/quinte: due ore CONSECUTIVE di mate:
\begin{lstlisting}
forall c in classes where c.year >= 4:
    exists s1 in slots: exists s2 in slots:
        consecutive(s1, s2)
        and lesson(class==c.name, day==s1.day,
                   hour==s1.hour, subject==Matematica)
        and lesson(class==c.name, day==s2.day,
                   hour==s2.hour, subject==Matematica)
\end{lstlisting}

\item Linguistico anno 1: niente sabato:
\begin{lstlisting}
forall l in lessons
    where l.class.curriculum == Linguistico
      and l.day == 6: false
\end{lstlisting}

\item 1A: religione solo nelle prime 4 ore:
\begin{lstlisting}
forall l in lessons where l.class == 1A
                       and l.subject == Religione:
    l.hour <= 11
\end{lstlisting}
\end{enumerate}

\subsection{Vincoli aula (5 esempi)}

\begin{enumerate}
\item Lab Fisica 1: una sola lezione per slot:
\begin{lstlisting}
forall s in slots:
    count l in lessons
        where l.classroom == "Lab Fisica 1"
          and l.day == s.day and l.hour == s.hour:
        l <= 1
\end{lstlisting}

\item Matematica solo in aule taggate ``matematica'':
\begin{lstlisting}
forall l in lessons where l.subject == Matematica:
    "matematica" in l.classroom.tags
\end{lstlisting}

\item Storia delle terze: aula con proiettore:
\begin{lstlisting}
forall l in lessons where l.subject == Storia
                       and l.class.year >= 3:
    "proiettore" in l.classroom.tags
\end{lstlisting}

\item La palestra ospita SOLO Scienze motorie:
\begin{lstlisting}
forall l in lessons
    where l.classroom.type == "palestra":
    l.subject == Scienzemotorie
\end{lstlisting}

\item Lab chimica: max 2 lezioni concorrenti:
\begin{lstlisting}
forall s in slots:
    count l in lessons
        where l.classroom.type == "lab_chimica"
          and l.day == s.day and l.hour == s.hour:
        l <= 2
\end{lstlisting}
\end{enumerate}

\subsection{Vincoli materia (3 esempi)}

\begin{enumerate}
\item Mate distribuita: max 2 ore/giorno per classe:
\begin{lstlisting}
forall c in classes:
    forall d in days:
        count l in lessons
            where l.class == c.name
              and l.subject == Matematica
              and l.day == d.index: l <= 2
\end{lstlisting}

\item Mate preferibilmente entro la 5a ora:
\begin{lstlisting}
forall l in lessons where l.subject == Matematica:
    l.hour <= 12
\end{lstlisting}

\item Educazione fisica solo in palestra:
\begin{lstlisting}
forall l in lessons
    where l.subject == Educazionefisica:
    l.classroom.type == "palestra"
\end{lstlisting}
\end{enumerate}

\subsection{Vincoli gruppo / studenti (3 esempi)}

\begin{enumerate}
\item Studenti BES devono essere in un gruppo Sostegno (uso di
\texttt{exists g in s.groups}):
\begin{lstlisting}
forall s in students where "BES" in s.tags:
    exists g in s.groups: g.name == "Sostegno"
\end{lstlisting}

\item Studenti con debito mate quarto $\to$ gruppo Recupero:
\begin{lstlisting}
forall s in students
    where "debito_matematica_4" in s.tags:
    exists g in s.groups:
        g.name == "Recupero Matematica"
\end{lstlisting}

\item Studenti BES: nessuna lezione il sabato:
\begin{lstlisting}
forall l in lessons:
    forall s in students
        where l.class == s.class
          and "BES" in s.tags:
        l.day != 6
\end{lstlisting}
\end{enumerate}

\subsection{Vincoli relazionali / coupling (4 esempi)}

\begin{enumerate}
\item Rossi in Aula 12: solo Matematica:
\begin{lstlisting}
forall l in lessons where l.teacher == Rossi
                       and l.classroom == "Aula 12":
    l.subject == Matematica
\end{lstlisting}

\item Se Rossi in 3A allora deve usare lab fisica almeno una
volta:
\begin{lstlisting}
(exists l in lessons:
    l.teacher == Rossi and l.class == 3A)
=>
(exists l2 in lessons:
    l2.teacher == Rossi and l2.class == 3A
    and l2.classroom.type == "lab_fisica")
\end{lstlisting}

\item Inglese alle prime: max 3 ore/sett per (docente,classe):
\begin{lstlisting}
forall t in teachers where t.subject == "Inglese":
    forall c in classes where c.year == 1:
        count l in lessons
            where l.teacher == t
              and l.class == c.name: l <= 3
\end{lstlisting}

\item Religione alle prime: mai in 1a ora:
\begin{lstlisting}
forall l in lessons where l.subject == Religione
                       and l.class.year == 1:
    l.hour != 8
\end{lstlisting}
\end{enumerate}

\section{Errori comuni}

\begin{itemize}
\item \texttt{Atteso COLON, trovato OP (=)} -- usa \texttt{==}
nei confronti dopo \texttt{where} (\texttt{=} singolo e' lecito,
ma confonde il parser quando appare in contesti ambigui).
\item \texttt{sorgente sconosciuta 'foo'} -- sorgente non in
\texttt{\_VALID\_SOURCES} e non e' un path raggiungibile via
env. Controlla l'ortografia.
\item \texttt{oltre 1000000 atomi: probabile esplosione
combinatoria} -- quantificatori troppo annidati. Sposta filtri
nel \texttt{where} pi\`u esterno per ridurre il set; oppure
splitta il vincolo in pi\`u vincoli salvati separatamente.
\end{itemize}

Vedere \texttt{docs/general\_dsl.md} per la guida completa con
sezioni 1--12 (galleria di 30+ esempi commentati,
troubleshooting esteso, reference tabellare di tutti gli
attributi per ciascuna entit\`a).

\section{Compilazione DSL $\to$ CP-SAT (Step 1)}
\label{sec:dsl-cpsat}

A partire da Step 1 del piano multi-day, il DSL ha un
\emph{compilatore} verso CP-SAT
(\texttt{engine/dsl\_to\_cpsat.py}, classe
\texttt{DSLConstraintCompiler}). Mentre l'evaluator post-hoc
guarda una soluzione gi\`a costruita per dichiararla feasible o
violata, il compilatore aggiunge i vincoli \emph{durante} la
costruzione del modello CP-SAT, in modo che il solver eviti
strutturalmente le violazioni.

Il compilatore \`e usato in tre punti del sistema:
\begin{itemize}
\item dal \texttt{MonolithicSolver} (cap.~\ref{ch:metodo_cpsat})
  quando il flag \texttt{via\_dsl=True} \`e attivo;
\item dai pricer di Branch-and-Price
  (cap.~\ref{ch:tecniche_avanzate}) per applicare gli stessi
  vincoli al sotto-problema;
\item da \texttt{PhaseBDaySolver} quando si invoca con
  \texttt{via\_dsl=True} per riprodurre la pipeline legacy
  affiancata dal DSL.
\end{itemize}

Tutte le sorgenti delle lezioni gi\`a note (\texttt{teachers},
\texttt{classes}, \texttt{slots}, ecc.) sono accessibili al
compilatore allo stesso modo. La differenza chiave \`e che
quando una variabile slot \`e ancora indecisa, il compilatore
emette il vincolo CP-SAT corrispondente al posto di valutare
una verit\`a; quando un'espressione \`e \emph{completamente
statica} (per esempio \texttt{consecutive(s1, s2)} con \texttt{s1}
e \texttt{s2} legati staticamente nel \texttt{forall}) la
valutazione viene eseguita a tempo di compilazione e si emette
un vincolo solo se serve.

\section{Predicati di slot: \texttt{consecutive}, \texttt{same\_day}}
\label{sec:dsl-slot-preds}

Step 3a aggiunge due predicati built-in che il compilatore
risolve a tempo di compilazione quando entrambi gli argomenti
sono slot legati staticamente:

\begin{itemize}
\item \texttt{same\_day(s1, s2)} \`e vero se i due slot cadono
  nello stesso giorno della settimana.
\item \texttt{consecutive(s1, s2)} \`e vero se i due slot cadono
  nello stesso giorno e differiscono di un'ora ($|h_1-h_2|=1$).
\end{itemize}

Sono i mattoni con cui si esprimono ``ore consecutive'',
``stacco fra plessi'', ``coppia di ore quando ce ne sono due
nel giorno'', ``compresenza che deve avvenire nello stesso
giorno''. Gli argomenti accettati sono o tuple letterali
\texttt{(d, h)} oppure path puntati come \texttt{l.slot} che
restituiscono una \texttt{(day, hour)}.

Esempio: una classe deve avere mate in tre ore consecutive in
\emph{qualche} giorno della settimana:
\begin{lstlisting}
exists s1 in slots: exists s2 in slots: exists s3 in slots:
    consecutive(s1, s2) and consecutive(s2, s3)
    and lesson(class==3B, day==s1.day,
               hour==s1.hour, subject==Matematica)
    and lesson(class==3B, day==s2.day,
               hour==s2.hour, subject==Matematica)
    and lesson(class==3B, day==s3.day,
               hour==s3.hour, subject==Matematica)
\end{lstlisting}

\section{Path \texttt{l.classroom.plesso}}
\label{sec:dsl-classroom-plesso}

Ancora in Step 3a, l'evaluator post-hoc e il compilatore
risolvono il path chained \texttt{l.classroom.plesso}: per ogni
lezione \texttt{l}, l'aula assegnata viene letta dalla mappa
\texttt{classroom\_for\_slot} fornita da Phase B / Phase
classroom-assignment, e il plesso viene risolto via la tabella
\texttt{plessi\_data}. Senza queste due strutture il path
ritorna \texttt{None} e il compilatore emette un diagnostico.

L'esempio canonico \`e una commuting rule fra plessi (vedi
sez.~\ref{sec:plessi}): ``nessun docente pu\`o avere lezione in
Plesso B nell'ora successiva a una lezione in Plesso A'':
\begin{lstlisting}
forall l1 in lessons where l1.classroom.plesso == 1:
    forall l2 in lessons where l2.classroom.plesso == 2
        and l2.teacher == l1.teacher
        and consecutive(l1.slot, l2.slot):
        false
\end{lstlisting}

\section{Vocabolario di pattern canonici (Step 3b)}
\label{sec:dsl-canonici}

Step 3b introduce un \emph{vocabolario} di chiamate di funzione
che il compilatore riconosce in modo speciale ed emette
direttamente come gruppi compatti di vincoli CP-SAT. Ognuna di
queste forme corrisponde a un pattern frequente che, scritto in
DSL puro a colpi di \texttt{forall}, produrrebbe centinaia o
migliaia di clausole. La forma canonica \`e zero-drift rispetto
all'encoding hardcoded del solver legacy.

\begin{tabular}{lp{6.7cm}}
\toprule
\textbf{Pragma} & \textbf{Equivalente legacy} \\
\midrule
\texttt{no\_holes\_class("3B")} &
  catena non-crescente di booleani che vieta i buchi
  nell'orario settimanale della 3B (legacy
  \texttt{add\_class\_no\_holes}) \\

\texttt{class\_present\_at\_hour("3B", 11)} &
  se la 3B ha una lezione in un giorno, lo slot all'ora 11
  deve essere occupato (legacy HARD-3 / ``uscita dopo le 12'') \\

\texttt{class\_day\_load\_in("3B", 0, 4, 5, 6)} &
  in un giorno la 3B ha 0 ore (giorno libero) o tra 4 e 6 ore
  (legacy HARD-1+HARD-2) \\

\texttt{teacher\_max\_per\_day("Rossi", 5)} &
  Rossi ha al massimo 5 ore in un giorno
  (legacy \texttt{MAX\_PROF\_HOURS\_PER\_DAY}) \\

\texttt{cattedra\_max\_per\_day("Rossi", "3B", "Matematica", 2)} &
  la cattedra (Rossi, 3B, Mate) ha al massimo 2 ore in un
  giorno (legacy \texttt{MAX\_PER\_DAY\_TRIPLE}) \\

\texttt{subject\_pair\_must("3B", "Scienzemotorie")} &
  quando 3B ha 2 ore di motoria nello stesso giorno, devono
  essere consecutive (legacy HARD-B / \texttt{add\_motorie\_pair}) \\

\texttt{subject\_pair\_exists("3B", "Matematica")} &
  quando 3B ha 2+ ore di mate nello stesso giorno, almeno una
  coppia di consecutive deve esistere (legacy HARD-A /
  \texttt{add\_math\_italian\_pair}) \\
\bottomrule
\end{tabular}

\medskip

Inoltre il pragma \texttt{no\_holes\_all\_classes()} e
\texttt{class\_present\_at\_hour\_all(11)} agiscono in batch su
tutte le classi nel \emph{slot scope} corrente, utile da
incollare in cima alla lista di vincoli senza dover enumerare i
nomi.

\section{Plesso commuting via DSL (Step 3c)}
\label{sec:dsl-plesso-commuting}

Step 3c collega le \texttt{PlessoCommutingRule} al DSL: la
helper \texttt{plesso\_commuting\_rule\_to\_dsl(rule)} produce
una lista di clausole DSL equivalenti alla regola legacy. La
property zero-drift \`e verificata nella suite di test:
applicare la regola via DSL produce esattamente le stesse
combinazioni proibite della pipeline hardcoded
\texttt{plessi\_constraints}.

\section{Seed legacy HARD via DSL (Step 3d-3e)}
\label{sec:dsl-seed-legacy}

Step 3d-3e introduce \texttt{seed\_implicit\_hardcoded(profs)}
che, dato il dizionario delle cattedre, restituisce una lista
di stringhe DSL che rappresentano TUTTI i vincoli HARD
hardcoded del solver legacy. Combinato con il flag
\texttt{via\_dsl=True} di \texttt{MonolithicSolver} e
\texttt{PhaseBDaySolver}, permette di costruire un modello
CP-SAT completo \emph{esclusivamente dal DSL}: nessun ramo di
codice hardcoded \`e attivo. Le invarianti che si garantiscono:

\begin{itemize}
\item zero-drift sintattico: le tuple di slot proibite/forzate
  emesse dal path DSL coincidono bit-per-bit con quelle del
  path legacy;
\item zero-drift semantico: le soluzioni feasible accettate
  dai due path sono identiche, verificato sui benchmark
  \texttt{small}/\texttt{medium}/\texttt{huge};
\item determinismo: l'ordinamento delle stringhe DSL prodotte
  da \texttt{seed\_implicit\_hardcoded} \`e stabile, quindi
  l'ordine di inserimento dei vincoli nel modello CP-SAT \`e
  riproducibile run-to-run.
\end{itemize}

Questo \`e il \emph{vehicle} per la migrazione progressiva: nei
prossimi step il path hardcoded verr\`a deprecato a favore del
path DSL, e l'esistenza di \texttt{seed\_implicit\_hardcoded}
permette di farlo senza perdere la copertura dei vincoli.

\section{Stato dei vincoli SOFT}
\label{sec:dsl-soft-status}

Il DSL accetta sia regole HARD che regole SOFT (con peso
opzionale). A livello di \emph{evaluator post-hoc} il SOFT \`e
pienamente implementato: le violazioni di clausole SOFT
contribuiscono allo score globale della soluzione e sono
visibili nelle metriche restituite da \texttt{run\_metrics}.

A livello di \emph{compilatore CP-SAT} (Step 1+) il SOFT \`e
attualmente accettato dal parser ma il backend di compilazione
non emette ancora la traduzione a costo soft (penalit\`a
positiva sull'obiettivo). Il diagnostico
\texttt{soft constraint X: SOFT cost not yet wired} \`e emesso
dal compilatore per ogni regola SOFT incontrata; la regola
viene archiviata nel modello come \emph{label informativo}, ma
il solver non aggiusta l'obiettivo. Nei prossimi step la SOFT
verr\`a integrata: l'utente legge ora la guida tenendo presente
che il SOFT in DSL \`e semanticamente definito ma il suo
\emph{enforcement} via CP-SAT \`e ancora in corso.

In pratica: usare HARD per qualunque cosa che non possa essere
violata; per le SOFT, usare il DSL come ``contratto
dichiarativo'' che oggi la pipeline esegue post-hoc e domani
eseguir\`a anche durante la compilazione.

\section{Diagramma dell'architettura DSL}
\label{sec:dsl-archi}

\begin{verbatim}
+--------------------------------------------------+
|  UI (modal "+ Nuovo vincolo DSL")                 |
|  webui/frontend/.../GeneralConstraintsCard.svelte |
+--------------------+-----------------------------+
                     |
                     v
+--------------------+-----------------------------+
|  POST /api/constraints/general                   |
|  webui/backend/routers/constraints.py            |
+--------------------+-----------------------------+
                     |
                     v
+--------------------+-----------------------------+
|  Parser + AST                                    |
|  webui/backend/utils/general_dsl.py              |
+--------------------+-----------------------------+
                     |
                     +---------+---------+
                     v         v         v
              post-hoc    compiler    SOFT score
              evaluator   to CP-SAT   (post-hoc)
              (cattura    (Step 1+)
               violazioni
               HARD)
                     |         |
                     v         v
              +--------------------+
              |   CP-SAT solver    |
              |   (Mono / Pricer / |
              |    PhaseB / RF /   |
              |    Meta)           |
              +--------------------+
\end{verbatim}

Stesso AST, due backend: l'evaluator legge la soluzione e
risponde ``feasible/violato''; il compilatore aggiunge vincoli
al modello prima della solve. La grammatica \`e unica e i nomi
delle sorgenti sono identici nei due backend.

